\documentclass[12pt,a4paper]{article}

\usepackage{fontawesome5}
\providecommand{\faPencilAlt}{\faIcon{pencil-alt}}

\usepackage{styles/couleurs}
\usepackage{styles/mise_en_page}
\usepackage{styles/boites}
\usepackage{styles/maths}
\usepackage{styles/figures}

\usepackage{enumitem}
\usepackage{array}
\usepackage{tabularx}
\usepackage{amsmath}
\usepackage{amssymb}
\usepackage{graphicx}
\usepackage{tikz}
\usepackage{xparse}

\usetikzlibrary{calc,angles,quotes,arrows.meta}

\setlist[itemize]{leftmargin=1.4em,itemsep=0.15em,topsep=0.2em}
\setlist[enumerate]{leftmargin=*,itemsep=0.22em,topsep=0.25em}

\newcommand{\lignepointillee}{\noindent\dotfill\par\vspace{0.45em}}

\newcommand{\lignesreponse}[1]{%
\begin{tcolorbox}[
enhanced,colback=white,colframe=NeutreC!40,boxrule=0.5pt,
arc=1.5mm,left=3mm,right=3mm,top=2mm,bottom=1mm,
before skip=2pt,after skip=2pt]
\foreach \i in {1,...,#1}{\lignepointillee}
\end{tcolorbox}}

\newcommand{\titreseq}[1]{%
\begin{center}
{\sffamily\LARGE\bfseries\color{TitreC}Séquence 6 -- Cosinus d'un angle aigu}\\[0.2em]
{\sffamily\large\bfseries\color{SavoirC}#1}
\end{center}}

\begin{document}

\titreseq{Progression professeur -- Calculer des longueurs seulement}

\remarque{Durée indicative : sept séances d'apprentissage de 50 à 55 minutes,
deux interrogations courtes, une séance de consolidation et une évaluation finale.

La progression choisit une route volontairement progressive :
\textbf{le cosinus est d'abord compris comme un multiplicateur dépendant de l'angle},
puis seulement comme un rapport. Cette entrée limite les inversions de formule
et donne du sens au calcul du côté adjacent.}

\section*{Objectifs de fin de séquence}

L'élève doit pouvoir :

\begin{itemize}
  \item repérer un triangle rectangle dans une figure simple ou composée ;
  \item identifier l'hypoténuse ;
  \item identifier le côté adjacent à l'angle aigu étudié ;
  \item exprimer le cosinus d'un angle aigu comme le quotient
        « côté adjacent divisé par hypoténuse » ;
  \item utiliser la calculatrice en mode degré ;
  \item calculer le côté adjacent connaissant l'hypoténuse et l'angle ;
  \item calculer l'hypoténuse connaissant le côté adjacent et l'angle ;
  \item choisir entre une multiplication et une division ;
  \item rédiger un calcul de longueur en nommant le triangle rectangle utilisé ;
  \item contrôler la vraisemblance d'un résultat ;
  \item résoudre un problème où le triangle rectangle doit être modélisé ;
  \item combiner, dans un problème de synthèse, cosinus et théorème de Pythagore.
\end{itemize}

\section*{Prérequis à réactiver}

\begin{itemize}
  \item reconnaître un triangle rectangle et son angle droit ;
  \item savoir que l'hypoténuse est le côté opposé à l'angle droit
        et le plus long côté du triangle rectangle ;
  \item lire et nommer un angle ;
  \item résoudre un calcul du type \(a=x\times b\) lorsque \(a\) ou \(x\) est inconnu ;
  \item utiliser les parenthèses et l'arrondi sur la calculatrice ;
  \item connaître le théorème de Pythagore pour le problème final.
\end{itemize}

\section*{Déroulé réaliste}

\renewcommand{\arraystretch}{1.38}
\begin{tabularx}{\textwidth}{
|>{\centering\arraybackslash}p{1.05cm}
|>{\raggedright\arraybackslash}p{4.2cm}
|>{\raggedright\arraybackslash}X
|>{\raggedright\arraybackslash}p{3.2cm}|}
\hline
\textbf{Séance} & \textbf{Support principal} & \textbf{Enjeu didactique} & \textbf{Production}\\
\hline
1 & Activité 1 : diagnostic et problème de la rampe & Réactiver le triangle rectangle ; faire constater que Pythagore ne suffit pas avec une longueur et un angle. & Besoin d'un nouvel outil.\\
\hline
2 & Activité 2 : familles de triangles & Mesurer plusieurs triangles ayant le même angle et conjecturer la constance du quotient adjacent/hypoténuse. & Tableau de mesures et conjecture.\\
\hline
3 & Activité 3 : vocabulaire et définition & Stabiliser la dépendance du côté adjacent à l'angle choisi ; introduire le cosinus et la touche calculatrice. & Interrogation 1 en début de séance.\\
\hline
4 & Activité 4 : calculer le côté adjacent & Utiliser le cosinus comme multiplicateur : \(\text{adjacent}=\text{hypoténuse}\times\cos(\text{angle})\). & Calculs directs et contrôle de vraisemblance.\\
\hline
5 & Activité 5 : calculer l'hypoténuse & Comprendre la division comme opération inverse ; éviter l'apprentissage mécanique d'une formule isolée. & Calculs d'hypoténuses.\\
\hline
6 & Activité 6 : laboratoire d'erreurs & Travailler la rédaction, le choix du triangle, le mode degré, l'arrondi et le sens de l'opération. & Interrogation 2 en début de séance.\\
\hline
7 & Activité 7 : problèmes de terrain & Modéliser des situations simples puis combiner cosinus et Pythagore dans un problème de hauteur. & Solutions rédigées.\\
\hline
8 & Consolidation différenciée & Automatismes, calculs gradués, choix de méthode et analyse d'erreurs. & Préparation à l'évaluation.\\
\hline
9 & Évaluation finale & Mobiliser vocabulaire, formule, calcul, rédaction et modélisation. & Évaluation sur 20.\\
\hline
\end{tabularx}

\section*{Automatismes}

Rituel de cinq minutes, quatre à six questions :

\begin{itemize}
  \item séances 1 à 3 : angle droit, hypoténuse, côté adjacent ;
  \item séances 3 à 5 : écrire un cosinus et utiliser la touche \(\cos\) ;
  \item séances 4 à 7 : décider « multiplier ou diviser » avant de calculer ;
  \item séances 5 à 8 : rédaction en trois lignes et contrôle du résultat ;
  \item après la séquence : une question de vocabulaire et un calcul court par semaine.
\end{itemize}

\section*{Obstacles anticipés et réponses}

\begin{itemize}
  \item \textbf{Le côté adjacent serait toujours le même.}
        Faire changer l'angle étudié sur un même triangle.
  \item \textbf{Le cosinus serait une longueur.}
        Répéter qu'il s'agit d'un nombre sans unité, compris entre \(0\) et \(1\)
        pour un angle aigu.
  \item \textbf{L'élève choisit multiplication ou division selon la taille attendue.}
        Revenir à la phrase opératoire :
        « l'adjacent est l'hypoténuse multipliée par le cosinus ».
  \item \textbf{La formule est écrite sans préciser le triangle rectangle.}
        Imposer d'abord la phrase :
        « Dans le triangle \ldots{} rectangle en \ldots{} ».
  \item \textbf{La calculatrice est en radians.}
        Installer un contrôle rapide :
        \(\cos 60^\circ=0{,}5\).
  \item \textbf{L'élève arrondit trop tôt.}
        Conserver la valeur affichée et n'arrondir qu'à la dernière ligne.
  \item \textbf{Le résultat est plus grand que l'hypoténuse.}
        Institutionnaliser le test de vraisemblance avant toute validation.
\end{itemize}

\section*{Choix assumés}

\begin{itemize}
  \item Aucun calcul d'angle n'est demandé dans cette séquence.
  \item La mémorisation d'un moyen mnémotechnique n'est pas l'entrée principale.
  \item La preuve complète par triangles semblables ou par Thalès n'est pas exigée ;
        l'investigation donne du sens au rapport constant.
  \item Les problèmes complexes arrivent seulement après la maîtrise des deux cas directs.
\end{itemize}

\end{document}
